% ============================================================
% Full body sections (except Abstract and Conclusion) for
% "Faithfulness Under Paraphrase: Auditing the Consistency of
% AI Explanations in Text Classifiers"
%
% Order: Introduction -> Related Work -> Methodology -> Results ->
% Discussion. Paste the whole file into the main .tex of the Overleaf
% arXiv template, after \maketitle/abstract and before the Conclusion.
% Uses only \cite{} keys already in references.bib -- no new bib
% entries needed. Assumes \usepackage{booktabs} (standard in most
% arXiv templates); if unavailable, replace \toprule/\midrule/\bottomrule
% with \hline.
% ============================================================

\section{Introduction}
\label{sec:introduction}

Text classifiers built on large pretrained language models are widely
deployed for tasks such as sentiment analysis, content moderation, and
intent detection, but their internal decision process is opaque to the
people who rely on their outputs. Post-hoc feature attribution methods
such as LIME \citep{ribeiro2016why}, SHAP \citep{lundberg2017unified},
and Integrated Gradients \citep{sundararajan2017axiomatic} address this
by assigning each input word a score meant to reflect how much it
contributed to the model's prediction, and are used in practice both to
debug models and to build end-user trust in individual predictions.

These methods are typically evaluated and applied one sentence at a
time: an explanation is computed for a given input, and that explanation
is trusted to describe \emph{that} input. Implicit in this practice is
an assumption that if a sentence is reworded without changing its
meaning or the model's prediction, the resulting explanation should
still identify essentially the same reasoning. This assumption is
rarely tested directly. Prior work has shown that attributions can be
fragile under small, deliberately adversarial perturbations
\citep{ghorbani2019interpretation,adebayo2018sanity,slack2020fooling,ivankay2022fooling,ivankay2022estimating},
but whether everyday, non-adversarial paraphrasing -- the kind of
variation a deployed system encounters from ordinary users -- already
destabilizes explanations has, to my knowledge, not been systematically
measured for text classifiers, a gap recent surveys of explainable AI
have also noted \citep{palikhe2025towards}.

This paper closes that gap by introducing Faithfulness-Under-Paraphrase
(FUP), a metric that compares the word attributions a classifier's
explanation assigns to a sentence against the attributions it assigns to
a meaning-and-prediction-preserving paraphrase of that same sentence. I
apply FUP to a DistilBERT sentiment classifier \citep{sanh2019distilbert}
over the SST-2 benchmark \citep{socher2013recursive}, comparing SHAP and
Integrated Gradients attributions across paraphrases generated by
backtranslation. The paper is organized around the following research
questions:

\begin{itemize}
  \item \textbf{Primary RQ}: Would feature attribution explanations for
  the same text classifier using techniques such as SHAP and Integrated
  Gradients be consistent even after the input text is paraphrased to
  mean the same thing and produce the same prediction as before? Or
  would there be significant shifts in the explanations?
  \item \textbf{RQ1}: How sensitive are the SHAP and Integrated Gradients
  attributions to changes in the original input that maintain the same
  meaning?
  \item \textbf{RQ2}: Are SHAP and Integrated Gradients different in
  terms of their robustness to paraphrasing?
  \item \textbf{RQ3}: What conditions lead to increased instability --
  in particular, does it scale with the degree of paraphrasing relative
  to the source sentence?
\end{itemize}

Across 296 correctly-classified SST-2 sentences and two independent
backtranslation pivots (575 surviving original/paraphrase pairs), I find
that both SHAP and Integrated Gradients are measurably but not
catastrophically destabilized by paraphrase, that the two methods differ
in \emph{how} they are destabilized rather than in \emph{whether} they
are, and that instability scales consistently with how much a paraphrase
diverges from the original sentence's surface form. The remainder of the
paper is organized as follows: Section~\ref{sec:related-work} reviews
prior work on post-hoc explanation methods, their fragility, and how
faithfulness has been defined and measured; Section~\ref{sec:methodology}
describes the FUP pipeline in detail; Section~\ref{sec:results} reports
the empirical findings; and Section~\ref{sec:discussion} interprets
these findings, relates them to prior work, and discusses limitations
and future work.

\section{Related Work}
\label{sec:related-work}

\subsection{Post-hoc Feature Attribution for Text Classifiers}

Post-hoc explanation methods assign an importance score to each input
feature without requiring access to or modification of the model's
architecture. LIME \citep{ribeiro2016why} approximates a model's local
behavior around an input with an interpretable surrogate model. SHAP
\citep{lundberg2017unified} unifies several earlier attribution schemes
under a single game-theoretic framework based on Shapley values, and its
Partition explainer variant, used in this paper, is designed
specifically for text and hierarchical data. Integrated Gradients
\citep{sundararajan2017axiomatic} takes a different approach entirely,
attributing importance by integrating the model's gradients along a path
from a baseline input to the actual input, and satisfies axioms (such as
sensitivity and implementation invariance) that many earlier
gradient-based methods do not. Toolkits such as Captum
\citep{kokhlikyan2020captum} have made these gradient-based methods
easy to apply to standard PyTorch models, including the transformer
architectures used in modern text classifiers. This paper uses SHAP and
Integrated Gradients as its two representative, widely used attribution
methods, one from each family.

\subsection{Fragility and Adversarial Manipulation of Explanations}

A separate line of work has questioned how much these attributions can
be trusted at all. Ghorbani et al.\ \citep{ghorbani2019interpretation}
show that imperceptible, adversarially chosen perturbations can
drastically change an explanation while leaving the model's prediction
unchanged, and Adebayo et al.\ \citep{adebayo2018sanity} show with
randomization-based sanity checks that some popular saliency methods are
insensitive to the model's learned parameters, undermining their
validity as explanations of the model at all. Slack et al.\
\citep{slack2020fooling} go further and construct classifiers that
deliberately fool LIME and SHAP into hiding biased behavior. Closest to
the present study, Ivankay et al.\
\citep{ivankay2022fooling,ivankay2022estimating} show that text
classifier explanations can be manipulated by small, adversarially
optimized changes to the input text. All of this work perturbs the
input \emph{against} the explainer, searching for a worst-case change.
This paper instead asks what happens under a perturbation with no such
adversarial objective: a paraphrase produced independently of the
classifier or explainer, whose only requirement is that it preserve
meaning and the model's prediction.

\subsection{Faithfulness: Definitions and Metrics}

"Faithfulness" is used inconsistently across the explainability
literature. Jacovi and Goldberg \citep{jacovi2020towards} argue that
faithfulness claims should be tied to an explicit, falsifiable
definition rather than asserted informally, and propose a set of
properties a faithfulness metric should satisfy. Yeh et al.\
\citep{yeh2019infidelity} operationalize a related notion, infidelity,
as an explanation's ability to predict how the model's output changes
under significant input perturbations, together with a complementary
sensitivity measure. Krishna et al.\ \citep{krishna2022disagreement}
document the "disagreement problem": different explanation methods
frequently disagree with each other on the same input and model, which
complicates any claim that a single method's output is "the"
explanation. Lyu et al.\ \citep{lyu2024towards} survey this broader
faithfulness literature in NLP and note that operational, quantitative
faithfulness metrics remain an active and unsettled area. FUP is an
operational faithfulness-adjacent metric in this tradition: rather than
perturbing the input to test whether the explanation predicts a change
in output (infidelity) or comparing multiple methods on a fixed input
(disagreement), it holds the input's meaning and the model's output
fixed and asks whether the explanation itself is stable. Manna and Sett
\citep{manna2024faithfulness} similarly connect faithfulness to
adversarial sensitivity; FUP asks the complementary, non-adversarial
question of whether attributions are already unstable under ordinary
rewording, before any adversarial search is applied.

\subsection{Paraphrase and Behavioral Invariance Testing}

Using paraphrase as a controlled, meaning-preserving perturbation has
precedent in NLP evaluation, though not for explanations specifically.
PAWS \citep{zhang2019paws} constructs paraphrase pairs by word
scrambling and controlled edits specifically to have high lexical
overlap while sometimes changing meaning, the inverse profile from the
paraphrases studied here, which are meaning-preserving but can have
comparatively low lexical overlap. CheckList
\citep{ribeiro2020beyond} introduces behavioral testing for NLP models,
including invariance tests that check whether a model's
\emph{prediction} is unchanged under label-preserving perturbations such
as paraphrase. This paper applies the same invariance logic one level
up the pipeline: instead of testing whether the prediction is invariant
under paraphrase (which I require as a precondition), I test whether the
\emph{explanation} of that prediction is invariant.

\subsection{Positioning of This Work}

Ehsan and Riedl \citep{ehsan2024explainable} argue that the XAI field
needs to revisit its evaluation practices as the models it explains grow
larger and more widely deployed, and recent surveys
\citep{lyu2024towards,palikhe2025towards} point to the absence of
systematic, empirical audits of explanation stability under everyday
input variation such as paraphrase, as opposed to worst-case adversarial
inputs. This paper contributes such an audit: an operational
faithfulness metric (FUP), applied at a moderate scale (296 sentences,
two independent paraphrase pivots, two attribution methods) to a widely
used, off-the-shelf text classifier, with all code and data released
for reproducibility.

% ============================================================
% METHODOLOGY / RESULTS / DISCUSSION
% (identical content to methodology_results_discussion.tex)
% ============================================================

\section{Methodology}
\label{sec:methodology}

\subsection{Task Setup}

I study whether post-hoc feature attribution explanations for a text
classifier remain stable when the input sentence is paraphrased in a way
that preserves both its meaning and the classifier's predicted label.
Concretely, I compare word-level attributions produced by SHAP
\citep{lundberg2017unified} and Integrated Gradients (IG)
\citep{sundararajan2017axiomatic} for an original sentence against the
attributions produced for a meaning-preserving paraphrase of that same
sentence, and quantify how much the two attribution profiles agree. I
refer to this comparison as Faithfulness-Under-Paraphrase (FUP). This
setup lets me address the three research questions from
Section~\ref{sec:introduction}: how sensitive each method is to
paraphrasing (RQ1), whether SHAP and IG differ in that sensitivity
(RQ2), and whether the degree of instability scales with how much the
paraphrase diverges from the original (RQ3).

\subsection{Classifier and Data}

I use \texttt{distilbert-base-uncased-finetuned-sst-2-english}, a
publicly available DistilBERT checkpoint \citep{sanh2019distilbert}
fine-tuned for binary sentiment classification, off the shelf. I
deliberately did not fine-tune my own classifier: the contribution of
this work is the faithfulness audit, not the classifier, and reusing a
citable, widely used checkpoint keeps that part of the pipeline
uncontroversial and reproducible.

Sentences come from the SST-2 validation split \citep{socher2013recursive},
accessed through the GLUE benchmark loader. I restrict the sample to
sentences of 6--25 words (long enough to paraphrase meaningfully, short
enough to keep SHAP's evaluation cost manageable) that the classifier
predicts correctly, and draw them with a fixed random seed for
reproducibility. This selection yielded 296 sentences for the run
reported below.

\subsection{Paraphrase Generation}

I generate paraphrases via round-trip (back-)translation: the English
sentence is translated into a pivot language and back into English using
MarianMT translation models, through two independent pivots, German and
French, so that each original sentence yields up to two paraphrases with
different rewording. Backtranslation was chosen over a dedicated
paraphrase model or hand-written rewrites because it is a cheap,
well-established way to obtain meaning-preserving rewordings without
introducing my own bias into how the paraphrases are constructed.

Unlike prior work that studies explanation robustness under
\emph{adversarially crafted} perturbations
\citep{ivankay2022fooling,ivankay2022estimating,slack2020fooling},
backtranslation is not optimized against the explainer or the
classifier in any way: it is a naturalistic source of paraphrases, which
makes FUP a test of stability under ordinary linguistic variation rather
than under a worst-case search.

For every candidate paraphrase, I re-run the classifier and discard the
pair if the predicted label changes. This guarantees that any difference
in attributions reflects the explanation method's response to a
reworded but semantically and predictively equivalent input, not a
difference in what is being explained.

\subsection{Attribution Methods}

\textbf{SHAP.} I wrap the classifier in a Hugging Face
\texttt{TextClassificationPipeline} and pass it to \texttt{shap.Explainer},
which automatically selects a Partition explainer with a word-level text
masker \citep{lundberg2017unified}. The Partition explainer recursively
splits the sentence rather than enumerating all word subsets, which
keeps evaluation tractable, and returns attributions already aggregated
to whole words.

\textbf{Integrated Gradients.} I use Captum's
\texttt{LayerIntegratedGradients} \citep{kokhlikyan2020captum},
attributing with respect to DistilBERT's word-embedding layer (gradients
require continuous inputs, not token IDs) \citep{sundararajan2017axiomatic}.
The baseline is an all-\texttt{[PAD]} sequence with the original
\texttt{[CLS]} and \texttt{[SEP]} tokens kept in place, representing "no
information," with 50 interpolation steps. Because IG operates at the
subword level, I merge \texttt{##}-prefixed WordPiece continuation tokens
back into whole words before comparing against SHAP's word-level output.

\subsection{Faithfulness-Under-Paraphrase Metrics}

A paraphrase does not guarantee a clean word-for-word alignment with the
original, so I cannot compare attribution vectors position by position.
I use two complementary scores instead:

\begin{itemize}
  \item \textbf{Top-5 salience retention}: of the original sentence's
  five most important words (by $|\text{attribution}|$), what fraction
  still appear among the paraphrase's own five most important words
  (case-insensitive match)? This is the headline FUP score.
  \item \textbf{Common-word rank correlation}: restricted to words that
  appear verbatim (case-insensitive) in both sentences, the Spearman
  correlation between their attribution ranks in the original and in the
  paraphrase. This asks whether words that survive the paraphrase keep a
  consistent relative importance, independent of whether they are in the
  top 5.
\end{itemize}

Both scores are computed separately for SHAP and for IG, for each
surviving original/paraphrase pair.

\subsection{Degree of Rephrasing (RQ3)}

To test whether instability scales with how much a paraphrase diverges
from the original, I compute two independent surface-form distance
measures for every pair: normalized Levenshtein edit distance
(character-level edit distance divided by the length of the longer
sentence) and sentence-level BLEU of the paraphrase against the
original.\footnote{Word overlap ratio (shared word types / original word
types) is also reported as a simpler, complementary diagnostic.} These
are used only as independent variables to correlate against the FUP
scores, never as inputs to the classifier or the explainers.

\subsection{Statistical Testing}

Because the FUP scores are bounded fractions rather than normally
distributed continuous values, I use the Wilcoxon signed-rank test, a
non-parametric test for paired samples, to compare SHAP's and IG's
scores across all surviving pairs for RQ2, and Pearson correlation
between each FUP score and each degree-of-rephrasing variable for RQ3.

\section{Results}
\label{sec:results}

\subsection{Sample and Survival Rate}

Of 296 correctly-classified sentences $\times$ 2 pivots (592 candidate
pairs), 575 pairs (97.1\%) preserved the classifier's predicted label
after backtranslation; 17 pairs (2.9\%) were discarded because the
paraphrase flipped the prediction. The classifier's mean confidence on
the retained originals was high (mean predicted-class probability
0.991), consistent with the correctness filter selecting comparatively
easy, unambiguous examples.

\subsection{RQ1: Are Attributions Stable Under Paraphrase?}

Table~\ref{tab:fup-results} summarizes the FUP scores over all 575
pairs. Both methods show partial but clearly imperfect stability: on
average, SHAP retains 61.8\% and IG retains 61.7\% of the original's
top-5 important words under paraphrase (i.e., roughly 3 of 5), while
common-word rank correlation is higher for both methods (SHAP 0.830, IG
0.795), indicating that words which survive the paraphrase largely keep
their relative importance ordering even when the top-5 \emph{set}
itself shifts. Neither score is close to the ceiling of 1.0 for either
method, so I find that both SHAP and IG explanations are measurably --
though not catastrophically -- sensitive to meaning-preserving
paraphrase.

\begin{table}[t]
\centering
\caption{FUP scores over all 575 original/paraphrase pairs (296 sentences, two pivots).}
\label{tab:fup-results}
\begin{tabular}{lccc}
\toprule
Metric & Mean & Median & Std.\ dev. \\
\midrule
Word overlap (diagnostic)      & 0.795 & 0.800 & 0.134 \\
Normalized edit distance       & 0.203 & 0.184 & 0.122 \\
BLEU                           & 0.487 & 0.474 & 0.222 \\
SHAP top-5 salience retention  & 0.618 & 0.600 & 0.249 \\
IG top-5 salience retention    & 0.617 & 0.600 & 0.258 \\
SHAP rank correlation          & 0.830 & 0.888 & 0.178 \\
IG rank correlation            & 0.795 & 0.858 & 0.216 \\
\bottomrule
\end{tabular}
\end{table}

\subsection{RQ2: Do SHAP and IG Differ in Robustness?}

The two methods are statistically indistinguishable on top-5 salience
retention (Wilcoxon signed-rank test, $p = 0.812$): neither method is
more likely than the other to keep the same top-5 words after
paraphrasing. On rank correlation, however, SHAP is significantly more
stable than IG (SHAP mean 0.830 vs.\ IG mean 0.795, Wilcoxon
$p < 0.001$). The two FUP scores therefore give a mixed answer to RQ2:
SHAP and IG are equally prone to reshuffling \emph{which} words rank in
the top 5, but among words that persist across the paraphrase, SHAP
preserves their relative importance more consistently than IG does.

\subsection{RQ3: Does Instability Scale with Degree of Rephrasing?}

Both degree-of-rephrasing variables correlate with all four FUP scores
in the expected direction: greater surface-form change is associated
with lower faithfulness. Normalized edit distance correlates negatively
with SHAP top-5 retention ($r = -0.582$), IG top-5 retention
($r=-0.515$), SHAP rank correlation ($r=-0.330$), and IG rank
correlation ($r=-0.405$); BLEU, a similarity score, correlates
positively with the same four scores ($r = 0.575$, $0.533$, $0.340$,
$0.395$, respectively). Word overlap shows a similar pattern
($r \approx 0.58$--$0.66$ with the two top-5 retention scores). All three
degree-of-rephrasing variables therefore converge on the same
conclusion: paraphrases that reword the sentence more heavily produce
less faithful explanations, for both attribution methods.

\subsection{Pilot-Scale vs.\ Full-Scale Comparison}

An earlier 15-example pilot run (28 pairs) suggested a much larger gap
in favor of IG (top-5 retention 0.66 for IG vs.\ 0.57 for SHAP). At the
full sample size this gap disappears on top-5 retention and reverses on
rank correlation. I report this discrepancy explicitly because it is
itself a methodological finding: conclusions about which explanation
method is "more robust" drawn from a small pilot ($n < 15$ sentences)
would have been misleading here, and I would caution against relying on
similarly small samples in future faithfulness studies.

\section{Discussion}
\label{sec:discussion}

\subsection{Interpretation}

The central finding is that meaning-preserving, non-adversarial
paraphrase measurably destabilizes post-hoc attributions from both SHAP
and Integrated Gradients, even though neither the meaning of the
sentence nor the classifier's decision changes. This extends prior
findings that explanations are fragile under small, deliberately
adversarial input perturbations
\citep{ghorbani2019interpretation,adebayo2018sanity,slack2020fooling,ivankay2022fooling,ivankay2022estimating}
to a setting with no adversary at all: an everyday linguistic
transformation like backtranslation is sufficient to shift roughly two
of the top five "important" words on average. This is practically
relevant because backtranslation-style paraphrase is exactly the kind of
variation a deployed system encounters from real users, rather than a
worst-case input constructed to break the explainer.

The SHAP/IG divergence on rank correlation, absent on top-5 retention,
is consistent with the broader observation that different explanation
methods can disagree with each other even when applied to the same
model and input \citep{krishna2022disagreement}: here they agree on
\emph{how much} the top-ranked words churn, but disagree on how
consistently the \emph{relative} importance of surviving words is
preserved. A plausible explanation is architectural: SHAP's Partition
explainer estimates attributions from coalitions over the sentence
structure, which may be comparatively insensitive to small changes in
exact wording or word position, whereas IG's embedding-space gradients
are computed pointwise along each input's own interpolation path and may
be more sensitive to shifts in tokenization, word order, or local
context introduced by the paraphrase. I present this as a plausible
mechanism rather than a confirmed cause, since isolating the exact
source of the difference would require controlled perturbation
experiments beyond the scope of this study.

The RQ3 results reinforce a limitation inherent to any faithfulness
metric built on paraphrase: the amount of measured instability is
entangled with the amount of surface-form change the paraphrase
introduces. Since edit distance, BLEU, and word overlap all correlate
with the FUP scores in a consistent direction, this is not an artifact
of any single distance measure, but a structural property of comparing
attributions across differently worded inputs. I discuss this further as
a limitation below rather than treating the raw FUP scores as a
context-free measure of "how faithful" an attribution method is in
general.

\subsection{Relation to Prior Work}

Jacovi and Goldberg \citep{jacovi2020towards} argue that faithfulness
claims should be tied to an explicit, falsifiable definition rather than
asserted informally; FUP follows that spirit by defining faithfulness
operationally as agreement between attributions across a
meaning-and-label-preserving transformation, in the same family as
infidelity-style perturbation metrics \citep{yeh2019infidelity} but using
natural paraphrase rather than synthetic input masking as the
perturbation. Relative to Manna and Sett's adversarial-sensitivity
framing of faithfulness \citep{manna2024faithfulness}, FUP asks a
narrower but complementary question: not whether an attribution can be
adversarially manipulated, but whether it is already unstable under
ordinary, non-adversarial rewording. Ehsan and Riedl
\citep{ehsan2024explainable} and recent surveys of explainability
methods \citep{lyu2024towards,palikhe2025towards} note the absence of
systematic empirical audits of exactly this kind for text classifiers;
this experiment, together with its 296-sentence, two-pivot sample, is a
step toward closing that gap, though at a considerably smaller scale
than a full benchmark study.

The paraphrase-as-invariance-test framing here is also conceptually
related to behavioral testing methodologies such as CheckList
\citep{ribeiro2020beyond}, which probes model \emph{predictions} for
invariance under label-preserving perturbations; this work applies the
same invariance logic to model \emph{explanations} rather than to
predictions.

\subsection{Limitations}

Several limitations affect how far these results generalize. First, the
abstract for this project proposed evaluating on both SST-2 and IMDB
\citep{maas2011learning}; the results reported here are SST-2 only, due
to the additional compute and time the IMDB extension would have
required within the project timeline. Second, only one classifier
(DistilBERT fine-tuned for SST-2 sentiment) and one task (binary
sentiment classification) are evaluated; whether the same pattern holds
for other architectures, tasks, or the longer, more varied documents in
IMDB is untested. Third, paraphrases are restricted to backtranslation
through two pivot languages; other paraphrase sources, such as
dedicated paraphrase generation models or adversarial paraphrase
datasets like PAWS \citep{zhang2019paws}, might produce different
stability patterns, particularly PAWS-style examples that preserve high
lexical overlap while changing meaning, which is the opposite failure
mode from the one studied here. Fourth, because the sampling procedure
keeps only correctly-classified sentences, and the label-preservation
filter implicitly keeps the paraphrase correct as well, this dataset
cannot speak to whether faithfulness differs between correct and
incorrect predictions, a comparison the original research questions
raised but which the current sampling design does not support. Finally,
as discussed under RQ3, the FUP scores are correlated with how much a
paraphrase changes the sentence's surface form, so they should be read
as sensitive to paraphrase \emph{style} as well as to genuine explanation
instability.

\subsection{Future Work}

Extending the evaluation to IMDB and to additional classifiers would
test whether the SHAP/IG comparison in Section~\ref{sec:results}
generalizes beyond one model and one dataset. Comparing backtranslation
against PAWS-style adversarial paraphrases \citep{zhang2019paws} within
the same FUP framework would separate "high divergence, meaning
preserved" instability (this study) from "low divergence, meaning
changed" instability, and clarify whether attribution methods are more
sensitive to lexical change or to semantic change per se. Finally,
because degree of rephrasing confounds the current FUP scores, a
regression-adjusted or matched-paraphrase-difficulty version of FUP
would let future work separate the two factors more cleanly.
